\begin{abstractCN}
% 硕士学位论文的中文摘要一般约500～800字。摘要内容应涉及本项科研工作的目的和意义、研究思想和方法、研究成果和结论
近年来，深度学习在诸多医疗影像分析任务中取得了长足进步，
并在手术阶段识别、解剖结构追踪等任务上取得了突破性成果。
然而，深度学习在医疗视觉任务上的成功往往建立在大规模、高质量的专家标注数据集基础之上。
在临床实际应用中，医疗视频数据的构建常面临影像对比度低、动作幅度微小、数据分布极度失衡等困难，
使得深度模型难以学习到有效的视觉表征以完成精准的动作理解任务。
因此，如何深度融合医学领域先验，实现更高效的医疗视频动作表征学习，对推动智慧医疗系统的实际落地具有重要意义。
然而，针对辅助诊断的吞咽造影视频和面向术中导航的青光眼手术视频两种典型视觉数据，高效的表征学习仍存在以下挑战：
（1）在吞咽造影视频微动作定位中，X 射线影像的低对比度及关键结构的模糊遮挡，导致微弱动作信号极易被静态背景噪声淹没，使得模型难以精准捕捉时序边界；
（2）在眼科手术视频在线识别中，手术阶段时长的显著失衡以及不同操作间极高的外观相似性，导致模型在实时推理时极易产生漏检与误判。
针对上述挑战，本文研究内容如下：

对于吞咽造影视频分析，本文针对时序微动作定位任务中解剖结构模糊及弱目标特征淹没的问题，对解剖结构校准机制进行探索。
提出了一种名为 \sexyname 的全新框架，通过引入基于临床先验的骨架热图模态，为模型提供鲁棒的几何结构参照，
从而引导模型学习抗干扰能力更强的视觉表征。
此外，为保证多模态特征的高效校准与增强，该方法利用跨模态特征校准模块 \crossblockshort 和通道级增强机制 \channelblockshort，
通过骨架特征引导空间聚焦并显式剥离冗余的静态背景信号，实现对细微动作信号的有效放大。
在 VFSS 基准数据集上的实验表明，相较于现有方法，所提出的方法在平均 mAP 指标上取得了 5.3\% 的性能提升，并在外部验证集中证实了其卓越的泛化能力。

% 这里只谈了显微手术，但是没有聚焦青光眼
对于青光眼手术导航，本文针对在线手术动作识别任务中时序逻辑混淆及数据分布不均的问题，
对手术规程逻辑协同机制进行探索，提出了一种名为 \sexytwo 的动作检测框架。该方法旨在利用手术流程的结构化先验，
通过动态逻辑演化模块 \graphmoduleshort，提升模型对手术状态跳转及动作持续时长的感知能力，
从而引导模型学习逻辑连贯性更强的场景表征。此外，还提出了类别边界加权损失 \lossshort，以促进长短时动作特征的均衡学习。
在构建的 MIGS 数据集及公开数据集 OphNet 上的实验表明，相较于现有方法，本文所提出的方法分别取得了 0.9\% 和 1.9\% 的在线识别准确率提升。
\end{abstractCN}
\keywordsCN{医疗视频理解；吞咽造影视频；时序微动作定位；手术视频；在线动作识别}